%% This file contains a number of tweaks that are typically applied to the main document.
%% They are not enabled by default, but can be enabled by uncommenting the relevant lines.


%%
%% Official support for 2-column teaser figure
%%
\usepackage{cuted} %< for strip environment

%%
%% For auto-referencing of figures (see fig/teaser.tex and \label{} therein) 
%%
\usepackage{currfile} %< for \currfilebase macro

%%
%% caption and (global) spacing overrides
%%
\usepackage{caption} %< for "\captionof" commands (teaser)
% \setlength{\abovecaptionskip}{.5em} % space above caption
% \setlength{\belowcaptionskip}{1em} % space below caption

%%
%% Inline annotations; for predefined colors, refer to "dvipsnames" in the xcolor package:
%% https://tinyurl.com/overleaf-colors
%%
\newcommand{\red}[1]{{\color{red}#1}}
\newcommand{\todo}[1]{{\color{red}#1}}
\newcommand{\TODO}[1]{\textbf{\color{red}[TODO: #1]}}
%%
%% Disable for camera ready / submission by uncommenting these lines:
%%
% \renewcommand{\TODO}[1]{}
% \renewcommand{\todo}[1]{#1}

%%
%% Work harder in optimizing text layout. Typically shrinks text by 1/6 of page, enable
%% it at the very end of the writing process, when you are just above the page limit.
%%
% \usepackage{microtype}

%%
%% Fine-tune paragraph spacing.
%%
% \renewcommand{\paragraph}[1]{\vspace{.5em}\noindent\textbf{#1.}}

%%
%% Globally adjusts space between figure and caption.
%%
% \setlength{\abovecaptionskip}{.5em}

%%
%% Allows "the use of \paper to refer to the project name"
%% with automatic management of space at the end of the word.
%%
% \usepackage{xspace}
% \newcommand{\paper}{ProjectName\xspace}

%%
%% Commonly used math definitions.
%%
% \DeclareMathOperator*{\argmin}{arg\,min}
% \DeclareMathOperator*{\argmax}{arg\,max}

%%
%% Tighten underline.
%%
% \usepackage{soul}
% \setuldepth{foobar}

%%
%% Reconfigure paragraph (tweaks spacing)
%%

%%
%% [Harvest] 한국어(XeLaTeX)와 마인드맵 표시
%%
\usepackage{kotex}
% 한글 글꼴: 나눔고딕(OFL, fonts/ 폴더에 동봉 — 로컬과 Overleaf에서 같은 결과). 기본 은바탕보다 화면 가독성이 좋다(사용자 요청 2026-09-24). 한글은 기울이지 않는다(기울임 자리도 바로 선 글꼴).
\setmainhangulfont{NanumGothic}[Path=fonts/, Extension=.ttf, UprightFont=*-Regular, BoldFont=*-Bold, ItalicFont=*-Regular, BoldItalicFont=*-Bold]
\setsanshangulfont{NanumGothic}[Path=fonts/, Extension=.ttf, UprightFont=*-Regular, BoldFont=*-Bold, ItalicFont=*-Regular, BoldItalicFont=*-Bold]
\usepackage{enumitem}
\setlist{nosep,leftmargin=1.2em}
\newcommand{\intent}[1]{\textcolor{red}{#1}}      % 빨간색 = 사용자가 의도하고 말한 것
\newcommand{\idea}[1]{\textbf{지금 생각:} #1}
\newcommand{\chg}[1]{\textbf{바뀐 점:} #1}
\newcommand{\dirn}[1]{\textbf{방향:} #1}
\newcommand{\decide}[1]{\textbf{[결정 필요]} #1}
